\ifthenelse{\equal{\interactive}{Y}}{
    % \hspace{0pt}
    % \vfill
    % {\Large\textbf{Eigenständigkeitserklärung}}

    % \enquote{Ich, \textbf{\vollername}, versichere hiermit, dass ich die Klausur \textbf{\veranstaltung} selbständig und nur unter Verwendung der erlaubten Hilfsmittel in der vorgegebenen Bearbeitungszeit bearbeiten werde. 

    % Insbesondere versichere ich, keine unerlaubte Hilfe anderer Personen in Anspruch zu nehmen und während der Klausur mit keiner anderen Person außer dem Prüfungsleiter/der Prüfungsleiterin zu kommunizieren. 

    % Mir ist bekannt, dass eine unwahre Erklärung \underline{rechtliche Folgen} haben und insbesondere dazu führen kann, dass die Klausur als nicht bestanden bewertet wird. Darüber hinaus ist mir bekannt, dass der Prüfer/die Prüferin in Verdachtsfällen innerhalb der Beurteilungsfrist mündliche Nachfragen zum Stoffgebiet vornehmen kann.}
    
    % %\vspace{2cm}
    % \begin{center}
    %     \underline{\large Ich, \textbf{\vollername}, bestätige hiermit die Eigenständigkeitserklärung:}

    %     \textbf{{\color{red} CHECKBOX MARKIEREN \myarrow \declerationbox}}
    % \end{center}
    % %\vspace{0.3cm}
     
    % \vfill
    % \hspace{0pt}

    % \newpage
    \hspace{0pt}
    \vfill
    \begin{center}
        \textbf{KLAUSUR in \veranstaltung}

        \textbf{Dauer: 90 MINUTEN}\\
		\textit{Einreichungsfrist: 10 Minuten nach Prüfungsende}
    \end{center}
    \vspace{-0.3cm}
    Die Klausur besteht insgesamt aus \pageref{LastPage} Seiten. 
    Bei Single- und Multiple-Response-Teilaufgaben müssen Sie eine oder mehrere richtige Aussagen markieren. Die genauen Modalitäten zur Bepunktung werden bei jeder Aufgabe deutlich erklärt. 

    Die Klausur wird in deutscher Sprache und einer zusätzlichen englischen Hilfestellung ausgegeben. Bitte benutzen Sie entweder das deutsche \underline{oder} das englische Exemplar. Rechtlich bindend sind die Aufgabenstellungen in deutscher Sprache, die englische Version dient lediglich als Hilfestellung. Sie können Ihre Lösungen in deutscher oder englischer Sprache formulieren.

	Für die Klausur sind \underline{eigene Notizen und von uns bereitgestellte Materialien} als Hilfsmittel zugelassen. Sie dürfen ein Wörterbuch verwenden. Die Nutzung weiterer Hilfsmittel ist untersagt. Die Bearbeitungszeit beträgt 90 Minuten. Danach haben Sie 10 Minuten, um Ihre Lösungen mithilfe des unten genannten Formulares einzureichen.
    \hspace{0pt}

	\ifthenelse{\equal{\uploadurl}{}}{}{
        \vspace*{\fill}
        \fbox{\parbox{\textwidth}{%
            \begin{center}
				{\color{red}\textbf{\Large IHRE UPLOAD-URL - NICHT WEITERGEBEN}}
                \vspace*{10mm}

                {\large\url{https://vs.uni-due.de/uploadexam/\uploadurl}}

				\vspace*{10mm}
				{Laden Sie diese PDF im Formular unter dieser Adresse hoch, wenn Sie mit der Bearbeitung fertig sind. Weitere Inforamtionen finden Sie in den Informationsmaterialien im Moodle-Kurs.}
            \end{center}
        }}
        \vspace*{\fill}
        \newpage
    }
    \vfill
}
{
    \begin{center}
        \textbf{KLAUSUR in \veranstaltung}
    
        \textbf{Dauer: 90 MINUTEN}
    \end{center}
    \vspace{-0.3cm}
    \ifthenelse{\equal{\vollername}{}}{%
    Bitte füllen Sie vor Beginn das Deckblatt aus und überprüfen Sie die Vollständigkeit Ihrer Aufgabenblätter!}{%
    Bitte prüfen Sie vor Beginn auf dem Deckblatt, ob es sich um Ihre Klausur handelt und überprüfen Sie die Vollständigkeit Ihrer Aufgabenblätter!}
    Die Klausur besteht insgesamt aus \pageref{LastPage} Seiten. Am Ende der Aufgabenblätter sind leere Seiten angefügt. Diese können Sie als Konzeptpapier für ihre Notizen verwenden. Sollte der Platz bei einer Aufgabe nicht für Ihre Lösung reichen, so können Sie die Lösung auf diese Blätter auslagern. Kennzeichnen Sie in diesem Fall dies bei der Aufgabenstellung und machen Sie bei den Lösungen deutlich, auf welche Aufgabe Sie sich beziehen. Wenn Lösungen auf Zusatzblättern nicht eindeutig zugeordnet werden können, werden sie nicht bewertet.

    Bei Single- und Multiple-Response-Teilaufgaben müssen Sie eine oder mehrere richtige Aussagen markieren. Bei einer falsch markierten Aussage wird die Teilaufgabe mit 0 Punkten bewertet, unabhängig von etwaigen richtig markierten Aussagen. Markieren Sie Ihre Antwort bitte mit einem Kreuz. Sollten Sie Ihre Antwort korrigieren wollen, streichen Sie die gesamte Kästchenspalte \underline{deutlich} durch und zeichnen eine neue Kästchenspalte.
    \begin{center}
        \includegraphics[width=0.43\textwidth]{kreuze.png}
    \end{center}
    % (CB) Print only if both versions are provided
    \ifthenelse{\equal{\englishandgerman}{yes}}
    {Die Klausur wird in deutscher Sprache und einer zusätzlichen englischen Hilfestellung ausgegeben. Bitte benutzen Sie entweder das deutsche \underline{oder} das englische Exemplar. Rechtlich bindend sind die Aufgabenstellungen in deutscher Sprache, die englische Version dient lediglich als Hilfestellung. Sie können Ihre Lösungen in deutscher oder englischer Sprache formulieren.}
    {}

    Für die Klausur sind \underline{keine} Hilfsmittel zugelassen. Das Mitführen von portablen Kommunikationsgeräten jedweder Art ist nicht erlaubt. Sie dürfen ein Wörterbuch verwenden. Die Bearbeitungszeit beträgt 90 Minuten.

    Arbeiten in unleserlicher Schrift können nicht beurteilt werden. Die Lösungen dürfen nur mit Tinte oder dokumentenechtem Kugelschreiber geschrieben werden. Lassen Sie am Ende der Klausur alle Unterlagen und Blätter, ob beschrieben oder nicht, auf ihrem Tisch. Nehmen Sie keine Unterlagen aus dem Raum mit. Bei Täuschungs- oder Störversuchen wird der entsprechende Paragraph der Prüfungsordnung angewendet.
}


